% Source: https://artofproblemsolving.com/community/c3h1630590_largest_area
% Stable identity: aops-topic-1630590-q01
\textbf{(a)} A circle of circumference $L$ has radius $r=\dfrac{L}{2\pi}$ and area $\pi r^2=\dfrac{L^2}{4\pi}$.

\textbf{(b)} The regular $n$-gon has side $s=\dfrac {L}{n}$. Joining the centre to the vertices splits it into $n$ congruent isosceles triangles with apex angle $\dfrac{360^\circ}{n}$; the height of each (the apothem $a$) bisects that angle and the base, so
\[
a=\dfrac{s}{2\tan\dfrac{180^\circ}{n}} .
\]
The area is $n\cdot\dfrac{1}{2} sa=\dfrac{1}{2}La=\dfrac{L^2}{4n\tan\dfrac{180^\circ}{n}}$.

\begin{asy}
usepackage("amsmath");
size(8cm);
defaultpen(fontsize(10pt));
real L=2*pi;
draw(circle((0,0),1),linewidth(0.9));dot((0,0));draw((0,0)--(1,0),dashed);label("$r$",(0.5,0),N);
pair c=(3.0,0);real s=L/6;real R=s;
pair[] H;for(int i=0;i<6;++i)H.push(c+R*dir(60*i));
draw(H[0]--H[1]--H[2]--H[3]--H[4]--H[5]--cycle,blue+linewidth(0.9));
draw(c--(H[0]+H[1])/2,blue+dashed);label("$a$",c+0.45*dir(30),NW,blue);
dot(c,blue);label("$s$",(H[4]+H[5])/2,S,blue);
label("circle: $L=2\pi r$",(0,-1.3));label("hexagon: $L=6s$",c+(0,-1.45),blue);
label("\mbox{}",(min(currentpicture,user=true)+max(currentpicture,user=true))/2,black); // keeps the text after the figure black when LaTeX typesets the labels inline
\end{asy}

\textbf{(c)} Write $\theta=\dfrac{\pi}{n}$ (in radians), so $0<\theta\le\dfrac{\pi}{3}$. The polygon's area is smaller than the circle's exactly when $\dfrac{L^2}{4n\tan\theta}<\dfrac{L^2}{4\pi}$, that is, when $n\tan\theta>\pi=n\theta$, that is, when
\[
\tan\theta>\theta\qquad\text{for }0<\theta<\dfrac{\pi}{2}.
\]
This holds: in a circle of radius $1$, take the sector $OAB$ with angle $\theta$ at the centre $O$, and let the tangent at $A$ meet line $OB$ at $T$. The triangle $OAT$ has legs $OA=1$ and $AT=\tan\theta$ and contains the sector, so
\[
\dfrac{1}{2}\tan\theta=[OAT]>[\text{sector }OAB]=\dfrac{1}{2}\theta .
\]
Hence every regular $n$-gon of perimeter $L$ encloses less than $\dfrac{L^2}{4\pi}$, the area of the circle. $\blacksquare$

(For example, the square encloses $\dfrac{L^2}{16}\approx0.0625L^2$, the regular hexagon $\dfrac{L^2}{8\sqrt3}\approx0.0722L^2$ and the circle $\dfrac{L^2}{4\pi}\approx0.0796L^2$. The poster's guess is right in general: the \emph{isoperimetric inequality} says that every closed curve of length $L$ encloses area at most $\dfrac{L^2}{4\pi}$, with equality only for the circle; its proof for arbitrary curves is beyond this exercise.)
